\documentclass[amsmath,twocolumn]{aastex702}

\usepackage[nonumberlist,nogroupskip]{glossaries}
\makenoidxglossaries

\newglossaryentry{isolated}{
  name={isolated},
  description={Not interacting with a companion. Most known black holes are in binaries, so they are brighter and easier to find}
}

\newglossaryentry{Roman}{
  name={Roman Space Telescope},
  description={Roman space telescope, launched recently, pretty much a MUCH better hubble. See \url{https://en.wikipedia.org/wiki/Nancy_Grace_Roman_Space_Telescope}}
}

\newglossaryentry{GaiaDR4}{
  name={Gaia DR4},
  description={sorta a wacky one, this is a data release from the \href{https://en.wikipedia.org/wiki/Gaia_(spacecraft)}{Gaia mission}. It is a massive upcoming data release containing a ludicrous amount of data on various stars over time}
}


\newglossaryentry{LIGO}{
  name={LIGO},
  description={Laser Interferomater Gravitational Wave Observatory}
}

\begin{document}

\onecolumngrid
\thispagestyle{empty}
\begin{center}
{\large Author's Note}
\end{center}

\vspace{1em}

Firstly, thank you so much for taking the time to check this out :), I really appreciate it. If you have any questions at all, please reach out, always love talking about research.

This serves as a simplified, easier and shorter to read version of the full paper(\cite{kiato2026uncoveringphysicaloriginsblack})\footnote{this feels insane to be able to do}, which can be found at \url{https://arxiv.org/pdf/2610.07152}. I wanted to make this both as an excuse to do academic writing, but for anyone interested in the research I've done to actually be able to understand and read. Additionally, I decided to format and organize this in a very form as the full paper, in the hopes that one can use this either as a supplement/pre-reading for the full thing, if they're interested in attempting the full paper.

Also, due to the nature of astrophysics, there will (inevitably) be a lot of jargon. To help combat this, there will a glossary at the end that has definitions and/or links to further information. Additionally, some larger topics that I would consider 'precursor info' (such as binary stars) will have dedicated sections attached at the end. These sections will be hyperlinked, meaning that (if your pdf viewer supports it) you should be able to just click said topic and it will bring you to that section of the text. This is mostly to avoid tangents that distract from the currently discussed topic.

From this point onward, it will be formatted as an academic paper, and content will be (mostly) synced up with the paper itself.

\clearpage

\title{Uncovering the Physical Origins of Black Holes in the Milky Way for Astronomical Surveys, \textbf{Layman Version}}

\author[orcid=0009-0007-5979-7930,sname='North America']{Pierson Lipschultz}
\affiliation{Center for Interdisciplinary Exploration and Research in Astrophysics (CIERA), Northwestern University, 1800 Sherman Ave,
Evanston, IL 60201, USA}
\affiliation{Department of Astronomy, University of Illinois Urbana-Champaign, 1002 W. Green Street, Urbana, IL 61801}
\email[show]{piersonlipschultz@gmail.com}

\begin{abstract}

  In this paper we present a number of measured properties from in-depth simulations of \gls{isolated} black holes (BHs). 
  These BHs are especially relevant currently, as both \gls{Roman} and \gls{GaiaDR4} are very likely to detect/reveal a large quantity of ones previously undiscovered. The location of these black holes will be discovered by these telescopes but as equally important is the lives which they've lived. I.e., what (and where!) were they before becoming black holes.
  This is something that is incredibly hard to do through observational means (i.e. Observing the black hole with some sort of telescope), as (most) black holes become incredibly sedentary after becoming formed, and besides some basic things like mass, there not much one can really glean from them. 
This means that its nearly impossible to`work backwards` when it comes to a figuring out how these BHs formed, and instead it must be approached from the front.
The main way to do this is to use massive very powerful simulation tools to simulate massive quantities of candidates, and then evolve said candidates and see which ones end up being what we're looking for. It's sort of like throwing a metaphorical balls into our stellar evolution plinko device, and seeing which ones land in the spot we're looking for, and then looking at how our 10-15 balls ended up there. 

\end{abstract}

\section{Introduction}

It is thoroughly predicted that there are millions of isolated
BHs in our milkyway, with estimates being around 200 million, however, it is well known that these black holes were not always isolated, and in fact the vast majority of them had to had formed in binary systems (see Appendix~\ref{sec:binary-stars}).
This means that the quantity and quantities of BHs have a distinct reliance on the evolution and processes involved in binary systems, and a careful treatment must be applied it in order to ensure that accurate predictions for BH populations are made.
Additionally, to date, only 2 dozen black holes with masses similar to stars have had confirmed mass measurements from direct dynamic observations\footnote{i.e. looking at the system itself with some sort of telecope and measuring the orbits and motion of the BH + companion}, with the majority of these being found in \ref{sec:XrBs} X-ray binaries, where the black hole is releasing a massive amount of energy and light from siphoning off material from a companion star.
These X-ray binaries served as our main method of observation and measurement of black holes until LIGO in 2015, where we suddenly measured many more black holes and their masses.

However, both of these methods have their own problems and skewing, and recently efforts to find and analyze black holes has been directed towards a method called gravitational lensing (see \ref{sec:grav-lensing}), 
a process where by a black hole will briefly pass infront of another star, and in the process, greatly increase the brightness of said star. 
This has been applied best to the gaia mission, as from the data three non-interacting\footnote{i.e. are not taking mass on from their companion} have been discovered.
This process also lead to the discovery of one isolated black hole by Hubble.

\pagebreak

\printnoidxglossary[title={Glossary},sort=word]

\appendix

\section{Binary Stars} \label{sec:binary-stars}
\subsection{X-ray binaries} \label{sec:XrBs}
\section{Gravitational lensing} \label{sec:grav-lensing}

%% Precursor background section. Expand this with the layman explanation of
%% binary systems, since most isolated BHs originally formed in binaries.

\bibliography{references}{}
\bibliographystyle{aasjournalv7}

\end{document}
